% =====================================================================
%  Sara Reyes -- Academic CV (long-form)
%  Template: "academic"  |  Compile with pdfLaTeX on Overleaf
%  Self-contained: only this file is needed.
% =====================================================================
\documentclass[11pt,letterpaper]{article}

\usepackage[T1]{fontenc}
\usepackage[utf8]{inputenc}
\usepackage{charter}              % classic serif body font
\usepackage[margin=0.9in]{geometry}
\usepackage{titlesec}
\usepackage{enumitem}
\usepackage{xcolor}
\usepackage[hidelinks]{hyperref}
\usepackage{fancyhdr}
\usepackage{lastpage}

\definecolor{rule}{gray}{0.35}

% ---- Page style: name + page number in footer ----
\pagestyle{fancy}
\fancyhf{}
\renewcommand{\headrulewidth}{0pt}
\fancyfoot[L]{\small\textit{Sara Reyes -- Curriculum Vitae}}
\fancyfoot[R]{\small Page \thepage\ of \pageref{LastPage}}

% ---- Section style: small caps with a full-width rule ----
\titleformat{\section}{\large\scshape\raggedright}{}{0em}{}[{\color{rule}\titlerule}]
\titlespacing*{\section}{0pt}{12pt}{5pt}
\titleformat{\subsection}{\normalsize\itshape}{}{0em}{}
\titlespacing*{\subsection}{0pt}{8pt}{3pt}

\setlist[itemize]{leftmargin=1.4em, itemsep=1pt, topsep=2pt, parsep=0pt}
\setlength{\parindent}{0pt}
\setlength{\emergencystretch}{1.5em}

% ---- Entry macros ----
% \cventry{Title}{Dates}{Organisation}{Extra line (optional, may be empty)}
\newcommand{\cventry}[4]{%
  \textbf{#1}\hfill #2\\
  #3\ifx&#4&\else\\ \textit{#4}\fi\par\vspace{2pt}}
% Numbered publication list
\newenvironment{publist}[1]{\begin{enumerate}[label={[#1\arabic*]}, leftmargin=2.6em, itemsep=4pt]}{\end{enumerate}}
\newcommand{\me}[1]{\textbf{#1}}   % highlight your own name in author lists

\begin{document}

% ============================ HEADER ============================
\begin{center}
  {\LARGE\bfseries Sara Reyes}\\[4pt]
  Ph.D. Student, Electrical \& Computer Engineering, University of Washington\\[2pt]
  \href{mailto:reyess4@uw.edu}{reyess4@uw.edu} \quad$\cdot$\quad 206-250-6131 \quad$\cdot$\quad Seattle, WA
\end{center}

% ============================ EDUCATION =========================
\section{Education}
\cventry{Ph.D., Electrical Engineering}{Expected June 2027}
        {University of Washington, Seattle, WA}{}
\begin{itemize}
  \item Focus: Backscatter Communication and RFID
  \item Advisor: Prof.\ Matthew S.\ Reynolds
\end{itemize}

\cventry{B.S., Electrical Engineering}{June 2023}
        {University of Washington, Seattle, WA}{}
\begin{itemize}
  \item Focus: Communications
  \item Ronald E.\ McNair Scholar; STARS Scholar
\end{itemize}

% ============================ RESEARCH ==========================
\section{Research Interests}
Communication systems for brain--computer interfaces and backscatter communication. Current
research uses the existing digital I/O pins of FPGAs as backscatter modulators: combined with a
continuous-wave carrier, standard Bluetooth chipsets can receive the resulting BLE packets,
enabling low-cost, low-power backscatter communication with no added components.

\section{Research Experience}
\cventry{Graduate Researcher}{Sep 2023 -- Present}
        {Electrical \& Computer Engineering Department, University of Washington}
        {Faculty Mentor: Prof.\ Matthew S.\ Reynolds}
\begin{itemize}
  \item Use existing digital I/O pins on FPGAs to enable UHF-band backscatter communication without additional components.
  \item Generate modulated square waves through FPGA digital pins using Verilog.
  \item Mentor an undergraduate student on implementing Verilog designs on FPGA and on project data collection.
\end{itemize}

\cventry{Undergraduate Researcher}{May 2019 -- Jun 2023}
        {Electrical \& Computer Engineering Department, University of Washington}
        {Faculty Mentor: Prof.\ Matthew S.\ Reynolds}
\begin{itemize}
  \item Assisted in developing next-generation ultra-low-power wireless biomedical devices.
  \item Built an automated RF measurement system (Arduino, MATLAB) to characterize channel path loss and bandwidth, improving the reliability of wireless neural recorders in dense multipath environments.
  \item Co-authored a paper in \textit{IEEE Transactions on Antennas and Propagation}.
  \item Designed a PIFA and supported functional and thermal testing of RF tags for aerospace and military applications under strict size and budget constraints; tags outperformed commercial products in five of seven thermal tests.
\end{itemize}

\cventry{REU Researcher}{Jun 2021 -- Aug 2021}
        {Center for Neurotechnology, University of Washington}{}
\begin{itemize}
  \item Characterized channel path loss and bandwidth of wireless neural recorders in dense multipath using MATLAB and a vector network analyzer (VNA).
  \item Gathered VNA data at the Washington National Primate Research Center.
  \item Defined research objectives with Prof.\ Matthew Reynolds and presented findings at a university-wide symposium.
\end{itemize}

\cventry{Undergraduate Researcher}{Apr 2019 -- Oct 2019}
        {Earth \& Space Sciences Department, University of Washington}
        {Faculty Mentor: Michael McCarthy}
\begin{itemize}
  \item Built an electrical interface to read a sounding rocket's accelerometer data.
  \item Created and revised the accelerometer interface PCB schematic in Eagle CAD.
\end{itemize}

% ============================ PUBLICATIONS ======================
\section{Publications}
\subsection{Journal Articles}
\begin{publist}{J}
  \item A.\ Oursland, O.\ Ianchenko, \me{S.\,M.\ Reyes}, R.\ Douglas, T.\,J.\ Moody, M.\ Lee, and M.\,S.\ Reynolds,
        ``Automated Optimization of Electronic Mode Stirring for Enhanced Backscatter Link Margin in Reverberant Cavities,''
        \textit{IEEE Journal of Radio Frequency Identification}, vol.~8, pp.~217--225, 2024.
        \href{https://doi.org/10.1109/JRFID.2024.3382687}{doi:10.1109/JRFID.2024.3382687}
  \item J.\,D.\ Rosenthal, A.\ Pike, \me{S.\ Reyes}, and M.\,S.\ Reynolds,
        ``Electronic Mode Stirring for Improved Backscatter Communication Link Margin in a Reverberant Cavity Animal Cage Environment,''
        \textit{IEEE Transactions on Antennas and Propagation}, vol.~70, no.~1, pp.~621--630, 2022.
        \href{https://doi.org/10.1109/TAP.2021.3102047}{doi:10.1109/TAP.2021.3102047}
\end{publist}

\subsection{Conference Papers}
\begin{publist}{C}
  \item K.\,J.\ Ho, \me{S.\,M.\ Reyes Lopez}, and M.\,S.\ Reynolds,
        ``Simultaneous Doppler Velocimetry and Backscatter Communication with a Low-Cost X-Band Radar Module,''
        in \textit{2025 IEEE SENSORS}, 2025, pp.~1--4.
        \href{https://doi.org/10.1109/SENSORS59705.2025.11331138}{doi:10.1109/SENSORS59705.2025.11331138}
  \item J.\,D.\ Rosenthal, \me{S.\,M.\ Reyes}, and M.\,S.\ Reynolds,
        ``An All-Digital, Fully Synthesizable Wireless Analog Voltage Telemetry Sensor Tag Using ZeroScatter Bluetooth Low Energy (BLE) Backscatter,''
        in \textit{2025 IEEE International Conference on RFID Technology and Applications (RFID-TA)}, 2025, pp.~1--4.
        \href{https://doi.org/10.1109/RFID-TA63091.2025.11265892}{doi:10.1109/RFID-TA63091.2025.11265892}
  \item \me{S.\,M.\ Reyes}, C.\ Fong, J.\,D.\ Rosenthal, and M.\,S.\ Reynolds,
        ``ZeroScatter BLE: All-Digital, Fully Synthesizable Bluetooth Low Energy (BLE) Backscatter Communication in Verilog RTL for FPGA-Based Devices,''
        in \textit{2024 IEEE International Conference on RFID (RFID)}, 2024, pp.~7--10.
        \href{https://doi.org/10.1109/RFID62091.2024.10582510}{doi:10.1109/RFID62091.2024.10582510}
  \item \me{S.\,M.\ Reyes}, A.\ Oursland, M.\ Lee, T.\,J.\ Moody, and M.\,S.\ Reynolds,
        ``Dual-Polarized Electronic Mode Stirring for Improved Backscatter Communication Link Margin in a Reverberant Cavity Animal Cage Environment,''
        in \textit{2023 IEEE International Conference on RFID (RFID)}, 2023, pp.~25--29.
        \href{https://doi.org/10.1109/RFID58307.2023.10178597}{doi:10.1109/RFID58307.2023.10178597}
\end{publist}

% ============================ PRESENTATIONS =====================
\section{Presentations}
\begin{publist}{P}
  \item \me{S.\,M.\ Reyes}, A.\ Oursland, M.\ Lee, T.\,J.\ Moody, and M.\,S.\ Reynolds,
        ``Dual-Polarized Electronic Mode Stirring for Improved Backscatter Communication Link Margin in a Reverberant Cavity Animal Cage Environment,''
        technical session, 2023 IEEE International Conference on RFID, Seattle, WA, Jun.\ 2023.
  \item \me{S.\ Reyes}, J.\,D.\ Rosenthal, and M.\,S.\ Reynolds,
        ``Electronic Mode Stirring Improves Wireless Communications in an Animal Cage Environment,''
        poster, National Diversity in STEM Conference, San Juan, Puerto Rico, Oct.\ 2022.
  \item \me{S.\ Reyes}, J.\,D.\ Rosenthal, and M.\,S.\ Reynolds,
        ``Electronic Mode Stirring Improves Wireless Communications in an Animal Cage Environment,''
        poster, Gabriel E.\ Gallardo Symposium, University of Washington, Apr.\ 2022.
  \item \me{S.\ Reyes}, J.\,D.\ Rosenthal, and M.\,S.\ Reynolds,
        ``Electronic Mode Stirring Improves Wireless Communications in an Animal Cage Environment,''
        poster, ECE Research Showcase, University of Washington, Mar.\ 2022.
  \item \me{S.\ Reyes},
        ``Electronic Mode Stirring Improves Wireless Communications in an Animal Cage Environment,''
        poster, WomXn at the Forefront of ECE Research, University of Washington, Nov.\ 2021.
  \item \me{S.\ Reyes},
        ``Electronic Mode Stirring Improves Wireless Communications in an Animal Cage Environment,''
        poster, Center for Neurotechnology, University of Washington, Aug.\ 2021.
  \item \me{S.\ Reyes}, J.\,D.\ Rosenthal, A.\ Pike, and M.\,S.\ Reynolds,
        ``Electronic Mode Stirring for Improved Backscatter Communication Link Margin in Reverberant Cavity Animal Cage Environments,''
        poster, IEEE International Conference on RFID (virtual), Sep.\ 2020.
\end{publist}

% ============================ AWARDS ============================
\section{Fellowships \& Awards}
\begin{tabular}{@{}p{0.78\textwidth}@{}r@{}}
  Leung Tsang Fellowship in Electromagnetics & 2025 \\
  CyberPOWDER Fellowship & 2024 -- 2025 \\
  Best Poster Award, IEEE RFID 2024 & 2024 \\
  Irene C.\ Peden ECE Fellowship & 2023 -- 2024 \\
  Ronald E.\ McNair Scholar & \\
  STARS Scholar & \\
\end{tabular}

% ============================ EMPLOYMENT ========================
\section{Professional \& Departmental Experience}
\cventry{Graduate Staff Assistant}{2024 -- Present}
        {Graduate Recruiting \& Admissions, Electrical \& Computer Engineering, University of Washington}{}
\begin{itemize}
  \item Coordinate M.S.\ and Ph.D.\ admissions with the faculty Graduate Program Coordinator and Graduate Program Advisor.
  \item Recruit and train the graduate admissions Triage Committee, assign reviewers to applications, and keep the review process on schedule.
  \item Represent UW ECE at national affinity-group conferences (NSBE, SHPE, SWE) and at local recruiting events, and follow up with prospective students.
  \item Run two applicant-support initiatives: Virtual Office Hours (AMA-style sessions with Ph.D.\ students and faculty) and the Graduate Applicant Support Program (GASP), which pairs applicants with volunteer graduate mentors for application feedback.
  \item Draft the annual admissions and recruitment budget and partner with the ECE communications team to promote recruiting events.
  \item Plan graduate community events with the ECE Events Team, manage event promotion and social media outreach, and use student feedback to improve programming.
  \item Provide peer mentorship that connects graduate students with departmental resources and opportunities; support information sessions, Graduate Visit Days, and orientation.
\end{itemize}

\cventry{Platform Engineering Intern}{Summer 2025}
        {Impinj, Seattle, WA}{}
\begin{itemize}
  \item Built an experimental setup for 3D localization of passive RAIN RFID tags with a handheld reader.
  \item Wrote Python tools for theoretical tag localization using synthetic data.
\end{itemize}

\cventry{Telecommunications Summer Intern}{Jun 2022 -- Aug 2022}
        {Hargis Engineers Inc., Seattle, WA}{}
\begin{itemize}
  \item Created CAD models of telecommunications systems for high-rise buildings in AutoCAD.
  \item Produced the initial telecommunications design for all 23 floors of a building, including port placement and IDF/MDF layouts.
  \item Worked across engineering disciplines to translate client needs into updated schematics.
\end{itemize}

% ============================ TEACHING ==========================
\section{Teaching}
\cventry{Teaching Assistant, Introduction to Signal Conditioning}{Jan 2024 -- Mar 2024}
        {Electrical \& Computer Engineering, University of Washington}{}
\begin{itemize}
  \item Led lab sessions using oscilloscopes, power supplies, multimeters, and function generators.
  \item Held office hours and problem sessions.
\end{itemize}

% ============================ LEADERSHIP ========================
\section{Leadership \& Service}
\cventry{Planning Committee, WomXn at the Forefront of ECE Research}{Aug 2023 -- Mar 2024}
        {University of Washington}{}
\begin{itemize}
  \item Recruited keynote speakers and industry panelists and helped select panel topics.
\end{itemize}

\cventry{DEI Committee, Electrical \& Computer Engineering}{Sep 2023 -- Jun 2024}
        {University of Washington}{}
\begin{itemize}
  \item Promoted equitable access to education and resources in ECE through initiatives such as the DEI Small Grants Program.
\end{itemize}

\cventry{A Vision for Electronic Literacy \& Access (AVELA)}{Apr 2023 -- Present}{Seattle, WA}{}
\begin{itemize}
  \item Taught middle school students basic circuit design and breadboarding at the William Grose Center.
  \item Trained class facilitators on coding curriculum for the Denny International Middle School coding club.
\end{itemize}

\cventry{Controls \& Power Lead, Washington Hyperloop}{Oct 2019 -- Apr 2021}{University of Washington}{}
\begin{itemize}
  \item Led a team of 13 engineers developing electronics and software for a pod in SpaceX's Hyperloop Competition V.
  \item Implemented power supplies for a linear induction motor, hydraulic pumps, and sensors; designed and assembled power and sensor PCBs.
  \item Presented Washington Hyperloop and STEM careers to middle and high school students (TAF at Saghalie, 2021).
\end{itemize}

% ============================ SKILLS ============================
\section{Technical Skills}
\textbf{Software:} Python, Verilog, MATLAB, Java, Ansys HFSS, Cadence Virtuoso, Altium, Eagle CAD, AutoCAD\\
\textbf{Hardware:} Vector network analyzers, signal generators, FPGAs, Arduino, circuit assembly, soldering

\end{document}
